\section*{Chapitre \ref{ch:ll:build-tick-to-trade-measured} --- Construction~: Tick-to-trade, mesurée}

\begin{solution}{exo:ll:build-tick-to-trade-measured:1}
Le carnet, la stratégie, le risque et la passerelle prennent $160 + 100 + 47 + 33 \approx \qty{340}{\nano\second}$ à leur médiane, soit moins du quart
des \qty{1.5}{\micro\second} cumulés par les étapes internes (décodage et file compris), et moins du dixième de la médiane interne
sur tous les ordres, qui inclut les attentes des ordres suivants.
\end{solution}

\begin{solution}{exo:ll:build-tick-to-trade-measured:2}
Une étape est la différence de deux relevés pris sur des threads différents, souvent sur des cœurs différents.
Le \omterm{def:ll:cpu-microarchitecture:tsc}{compteur d'horodatage} (TSC) est un compteur unique, invariant et synchronisé entre les cœurs de cette machine (chapitre~5), donc la différence est une durée.
Des horloges par thread, avec des décalages ou des cadences différents, feraient dériver les étapes, et certaines deviendraient négatives.
\end{solution}

\begin{solution}{exo:ll:build-tick-to-trade-measured:3}
Le $k$-ième événement du paquet est publié après $k \times \qty{100}{\nano\second}$, si bien que les cinq vont de 100 à 500, en moyenne
\qty{300}{\nano\second}. Côté moteur, les mêmes cinq événements sont traités à la suite, chacun attendant le traitement de ceux qui précèdent,
et l'envoi de leurs ordres.
\end{solution}

\begin{solution}{exo:ll:build-tick-to-trade-measured:4}
L'étape de risque, mesurée de la décision à la vérification du second ordre~: elle inclurait tout l'envoi du premier
ordre, soit environ \qty{4}{\micro\second}. Si la moitié des ordres sont dans ce cas, la médiane de l'étape tomberait entre les deux groupes,
quelque part entre les dizaines de nanosecondes de la seule vérification et les microsecondes de l'envoi~: l'attribution serait fausse, et non la barrière.
\end{solution}

\begin{solution}{exo:ll:build-tick-to-trade-measured:5}
En gardant environ \qty{100}{\nano\second} pour le décodage proprement dit~: $0.9 + 0.8 + (0.10 + 0.16 + 0.10 + 0.05 + 0.03) \approx
\qty{2.1}{\micro\second}$, dans le budget, avec environ \qty{0.9}{\micro\second} de marge.
\end{solution}

\begin{solution}{exo:ll:build-tick-to-trade-measured:6}
Les deux horodatages sont pris par logiciel autour des appels système~: celui de la bourse avant \texttt{sendto}, celui du thread du flux après
\texttt{recv}, si bien que l'étape contient l'appel de l'émetteur, le chemin de bouclage du noyau et l'appel du récepteur, sans aucun relevé
entre les deux. Un \omterm{def:ll:measuring-latency:hwts}{horodatage matériel}, pris par la carte quand le paquet est sur le câble, sépare la part de l'émetteur et
celle du récepteur, et sur un vrai réseau ajoute le câble lui-même.
\end{solution}

\begin{solution}{exo:ll:build-tick-to-trade-measured:7}
Un seul message de file par paquet, avec jusqu'à cinq événements~; le moteur les traite dans l'ordre, si bien que les ordres
ne changent pas (le hachage du test le vérifie). L'étape de décodage de tous les événements sauf le dernier disparaît, puisque le lot est publié d'un coup, et l'étape de file
perd le surcoût par message~; la file d'attente derrière les envois des ordres précédents subsiste.
\end{solution}

\begin{solution}{exo:ll:build-tick-to-trade-measured:8}
Elle ne mesure que la partie interne au processus, en laissant de côté la carte réseau, la réception et l'émission du noyau (sur ce portable, la plus grande part
du total) et le câble~; elle s'arrête à l'appel de \texttt{send}, et non au départ de l'ordre~; une moyenne ne dit rien de
la queue, là où les courses sont perdues. Il faut reporter les centiles de l'intervalle entre le paquet sur le câble et l'ordre sur le câble.
\end{solution}

\begin{solution}{pb:ll:build-tick-to-trade-measured:1}
\begin{enumerate}
\item Réception (de exch à recv), décodage (de recv à pub), file (de pub à pop), carnet (de pop à book), stratégie (de book à decide), risque (de decide
à risk), passerelle (de risk à encode), émission (de encode à venue).
\item Tick-to-trade~: environ \qty{8.8}{\micro\second}, \qty{79}{\micro\second} et \qty{3.5}{\milli\second}~; en interne~: environ
\qty{3.7}{\micro\second}, \qty{42}{\micro\second} et \qty{3.5}{\milli\second}.
\item L'émission a la médiane la plus élevée (environ \qty{2.9}{\micro\second})~; la file a le centile 99.9 le plus élevé (environ
\qty{3.5}{\milli\second}).
\item Quand le processeur du thread moteur est repris (le noyau, l'hyperviseur, une autre tâche), les événements attendent dans la file que le
thread reprenne la main~: c'est l'ordonnancement de la machine, non le code.
\item La stratégie et la \omterm{def:ll:pre-trade-risk-and-kill-switches:gate}{barrière de risque} atteignent leurs cibles~; le décodage manque la sienne d'un facteur d'environ 7, le carnet d'un tiers, la réception d'un facteur 1.6 et
l'émission d'un facteur 3.6.
\item Les étapes d'un ordre se composent exactement en son propre total~; ne retenir que les premiers ordres garde la composition cohérente avec
l'attribution par étape, tandis que la distribution du tick-to-trade inclut les attentes des seconds ordres.
\item Les quantiles d'une somme ne sont pas la somme des quantiles (chapitre~1)~: les cas lents de différentes étapes coïncident rarement, et quand
ils le font, ce ne sont pas les mêmes messages.
\item Qu'un chemin atteigne, ou n'atteigne pas, un objectif de performance annoncé sur un matériel annoncé, étape par étape~: une propriété du
système que nul test fonctionnel ne voit.
\item $0.9 + 0.8 + 1.5 \approx \qty{3.2}{\micro\second}$ à la médiane.
\item Environ les deux tiers des \qty{1.5}{\micro\second} en interne sont de l'attente (les événements du paquet qui s'attendent mutuellement dans
le décodage et sur la file)~; le reste est du calcul, dont les composantes du carnet prennent environ \qty{340}{\nano\second}.
\item Ils supprimeraient la préemption qui forme la queue de la file~; ils ne changeraient ni les médianes du calcul, ni la
réception et l'émission du noyau.
\item Il ajoute aux mesures des chiffres tirés d'un budget~; seule une exécution du harnais sur le serveur cible, avec sa carte et son
réglage, les mesure.
\item \textbf{Résultat nommé.} Sur ce portable, par ailleurs inactif, sur plus de 10\,500 ordres~: réception environ 1.4, décodage 0.44, file 0.76, carnet
0.16, stratégie 0.10, risque 0.05, passerelle 0.03 et émission \qty{2.9}{\micro\second} à la médiane~; tick-to-trade environ
\qty{8.8}{\micro\second} à la médiane, \qty{79}{\micro\second} au centile 99 et \qty{3.5}{\milli\second} au centile 99.9~;
projeté pour un serveur réglé avec kernel bypass, environ \qty{3.2}{\micro\second} à la médiane, \qty{2.1}{\micro\second} sans
la file du paquet.
\item Environ \qty{340}{\nano\second} de calcul par les composantes du carnet, contre quelque \qty{1.2}{\micro\second} de décodage et
d'attente dans les \qty{1.5}{\micro\second} des étapes du premier ordre~; la médiane interne de tous les ordres, \qty{3.7}{\micro\second}, est
surtout de l'attente.
\item Émission et réception~: \omterm{def:ll:linux-tuning:polling}{contournement du noyau} (kernel bypass, chapitre~13 et le livre~14)~; la queue de la file~: \omterm{def:ll:linux-tuning:isolation}{isolation de cœurs} (chapitre~13)~; la
file d'attente~: regroupement des transmissions (chapitres~12 et 18).
\item Sinon, la mesure pourrait porter sur un chemin qui calcule autre chose, ou qui saute du travail quand quelque chose tourne mal.
\item Une allocation sur le chemin critique est un danger pour la latence (chapitre~6), qu'une modification peut introduire en silence~; le compte
la détecte au passage suivant.
\item Les étapes de réception et d'émission avec \omterm{def:ll:measuring-latency:hwts}{horodatages matériels}, et les queues avec les cœurs du chemin isolés.
\item Les chiffres par composante des chapitres précédents~: chacun a été mesuré seul, avec des caches chauds et sans voisins.
\item Mesurer sur le chemin qui calcule la bonne réponse, étape par étape, avec une seule horloge et en centiles.
\end{enumerate}
\end{solution}

\begin{solution}{iq:ll:build-tick-to-trade-measured:1}
Le temps entre l'arrivée de l'événement de marché dans la firme (idéalement le paquet sur le câble, à la carte réseau) et
l'ordre qui en résulte, jusqu'à son départ de la firme (l'ordre sur le câble). On part et on arrête avec des \omterm{def:ll:measuring-latency:hwts}{horodatages matériels} quand c'est possible~; sinon on dit précisément quels
points logiciels sont utilisés et ce qu'ils laissent de côté.
\iqlookfor{du câble au câble, et honnêteté sur les points de mesure.}
\end{solution}

\begin{solution}{iq:ll:build-tick-to-trade-measured:2}
On lit une horloge commune et peu coûteuse (le \omterm{def:ll:cpu-microarchitecture:tsc}{compteur d'horodatage}) à quelques frontières d'étape, on porte les relevés avec le message dans
des enregistrements préalloués, on les écrit hors du chemin critique (un \omterm{def:ll:logging-capture-and-replay:binlog}{journal binaire}), puis on les joint par identité de message~; on attribue les étapes par message
et on publie des centiles.
\iqlookfor{peu de points, une horloge, enregistrement hors chemin, attribution par message.}
\end{solution}

\begin{solution}{iq:ll:build-tick-to-trade-measured:3}
Sur ce que le chemin attend plutôt que sur ce qu'il calcule~: préemption et interruptions sur ses cœurs (isolation), \omterm{def:ll:cpp-for-latency-i-memory:fault}{défauts de page} et
allocation, files d'attente derrière les rafales, pile réseau du noyau, \omterm{def:ll:managed-and-functional-languages-on-the-desk:gc}{ramasse-miettes} ou journalisation sur le chemin. L'étape dont la
queue s'allonge désigne laquelle.
\iqlookfor{ordonnancement et files d'attente avant le code.}
\end{solution}

\begin{solution}{iq:ll:build-tick-to-trade-measured:4}
Du câble à la carte et DMA, puis une réception par \omterm{def:ll:linux-tuning:polling}{contournement du noyau} (moins d'une microseconde)~; décodage de dizaines de nanosecondes~; carnet et stratégie de dizaines à
quelques centaines~; risque, dizaines~; encodage, dizaines~; envoi par \omterm{def:ll:linux-tuning:polling}{contournement du noyau} et de la carte au câble (moins d'une microseconde). Quelques microsecondes au total, dont l'essentiel
se joue dans le chemin réseau.
\iqlookfor{ordres de grandeur par étape, et où passe le temps.}
\end{solution}

\begin{solution}{iq:ll:build-tick-to-trade-measured:5}
À partir de l'exigence concurrentielle aux centiles qui comptent, répartie entre les étapes par la borne de l'union pour les queues
(chapitre~1)~; vérifiée par un harnais qui horodate chaque étape sur le chemin de production et contrôle chaque mise en production par rapport au budget.
\iqlookfor{cibles par étape et par centile, et un harnais qui les contrôle.}
\end{solution}

\begin{solution}{iq:ll:build-tick-to-trade-measured:6}
\omterm{def:ll:measuring-latency:hwts}{Horodatages matériels} à la carte pour le câble en entrée et en sortie~; \omterm{def:ll:cpu-microarchitecture:tsc}{compteur d'horodatage} (TSC) à chaque frontière d'étape, calibré sur une
horloge traçable en UTC~; horodatages portés avec le message et écrits dans un \omterm{def:ll:logging-capture-and-replay:binlog}{journal binaire} par un thread séparé~; jointure par numéro de séquence
et identifiant d'ordre~; histogrammes par étape et par mise en production, avec alertes sur les queues~; le même harnais dans la barrière de performance.
\iqlookfor{horloges, jointures, histogrammes, et lien avec le processus de mise en production.}
\end{solution}
